\documentclass[11pt,a4paper]{article}
\usepackage[T1]{fontenc}
\usepackage[utf8]{inputenc}
\usepackage{lmodern,microtype}
\usepackage[margin=26mm,headheight=14pt]{geometry}
\usepackage{amsmath,amssymb,amsthm,mathtools,mathrsfs}
\usepackage{booktabs,longtable,array,enumitem,needspace}
\usepackage{xcolor,xurl}
\definecolor{inkblue}{RGB}{25,54,85}
\usepackage[colorlinks=true,linkcolor=inkblue,citecolor=inkblue,urlcolor=inkblue]{hyperref}
\usepackage{fancyhdr}
\pagestyle{fancy}\fancyhf{}
\fancyhead[L]{\small Resonant polylogarithm inversion}
\fancyhead[R]{\small ProveIt research contribution}
\fancyfoot[C]{\thepage}
\setlength{\parskip}{4pt}\setlength{\parindent}{1.1em}
\setlist{itemsep=3pt,topsep=5pt}
\numberwithin{equation}{section}
\newtheorem{theorem}{Theorem}[section]
\newtheorem{proposition}[theorem]{Proposition}
\newtheorem{lemma}[theorem]{Lemma}
\newtheorem{corollary}[theorem]{Corollary}
\theoremstyle{definition}\newtheorem{definition}[theorem]{Definition}
\theoremstyle{remark}\newtheorem{remark}[theorem]{Remark}
\DeclareMathOperator{\Li}{Li}
\DeclareMathOperator{\Res}{Res}
\DeclareMathOperator{\Log}{Log}
\DeclareMathOperator{\Rea}{Re}
\newcommand{\C}{\mathbb C}\newcommand{\R}{\mathbb R}\newcommand{\N}{\mathbb N}
\newcommand{\dd}{\,\mathrm d}\newcommand{\eps}{\varepsilon}
\newcommand{\F}{\mathcal F}\newcommand{\B}{\mathcal B}
\newcommand{\A}{\mathcal A}\newcommand{\E}{\mathcal E}
\newcommand{\HH}{\mathcal H}
\newcommand{\coef}[2]{[#1]#2}
\newcommand{\Ot}{O_{\mathrm{tot}}}
\hypersetup{pdftitle={Resonant Polylogarithm Inversion: Uniform Order Crossover, Gamma-Zeta Fold Loci, and Convergent Exponential Sectors},pdfauthor={Research contribution prepared for Vladimir Reshetnikov with OpenAI assistance}}
\begin{document}
\begin{center}
{\small\scshape Analysis / polylogarithms / transseries}\par\vspace{10pt}
{\LARGE\bfseries Resonant Polylogarithm Inversion}\par\vspace{7pt}
{\Large Uniform Order Crossover, Gamma--Zeta Fold Loci,\par and Convergent Exponential Sectors}\par\vspace{14pt}
{\normalsize Prepared for Vladimir Reshetnikov with OpenAI assistance}\par
{\small 11 October 2026}\par
\end{center}
\begin{abstract}
The familiar expansion of a polylogarithm at its unit argument is not a
uniform inverse expansion when its order approaches an integer. We retain
the cancelling Gamma and zeta poles as one holomorphic block, derive every
integer-order logarithmic jet, and construct a convergent inverse chart in
three small variables. The chart's crossover function is
$\A(\lambda)=(1-e^{-\lambda})/\lambda$; its nonzero complex zeros require a
different scaling.

At these zeros we obtain an exact Gamma--zeta formula for the critical locus
of the singular core at every integer order, and compute the first
analytic-tail deformation of that locus. The dilogarithm is exceptional:
its undetuned secondary resonance is already a fold. If
$y=e^{-X}$, $\lambda_k=2\pi i k\ne0$, and $\eps=\lambda_k/X$, the equation
$\zeta(2+\eps)-\Li_{2+\eps}(e^{-t})=y$ has two roots near $t=y$, with
$t/y=1\pm\sqrt{-2K\lambda_k/X}+O(X^{-1})$. Here
$K=\gamma_1+((\gamma-1)^2+\zeta(2)+1)/2>0$.
We give the next coefficients, a convergent two-generator transseries,
the leading coefficient of every exponential sector, and a moving-fold
normal form with local branch interchange. Finite residue formulas determine
all chart coefficients from finite Taylor data. Proofs, executable exact
checks, independent floating-point comparisons, an audit of the inspected
repository material, and further research questions are included.
\end{abstract}
\noindent\textbf{Status.} The theorems are proved in this manuscript, but this is
an unrefereed research contribution, not a Lean formalization. Classical
continuation, Gamma expansions, implicit-function arguments, and analytic
fold mechanisms are credited rather than claimed as new. The model-specific
resonance classification, explicit critical-locus identities, and sector
formulas are the principal candidate contributions; publication priority has
not been established by an exhaustive literature search.
\clearpage
\begingroup
\small
\setlength{\parskip}{0pt}
\tableofcontents
\endgroup
\clearpage

\section{Problem, contribution, and relation to the repository}
\label{sec:intro}
The two selected repository areas provide complementary ingredients. The
polylogarithm project organizes exact identities, parameter derivatives,
Gamma and polygamma relations, and Stieltjes constants. The transseries
project develops support-controlled inversion, analytic realizations,
critical charts, and exponential sectors. Its critical-Hahn article already
uses polylogarithms with a fixed nonintegral order to build a convergent
multivariable inverse \cite{repo-poly,repo-trans,repo-hahn}. We do not claim
that polylogarithmic inversion or the analytic implicit-function method is
absent from the repository.

The question addressed here is narrower and different: what happens when
\emph{the order itself approaches an integer at the same time as the
argument approaches one}? Two effects must be kept separate. At an integer,
a Gamma pole cancels one zeta-series pole and produces logarithms. At certain
complex values of the scaled order, the leading logarithmic coefficient
cancels again. That second cancellation changes the inverse scale and, for
the dilogarithm, creates a two-sheet chart.

\subsection{A representative result}
Let $\gamma_1$ use the standard Laurent convention
\begin{equation}\label{eq:stieltjes}
 \zeta(1+z)=\frac1z+\sum_{r\ge0}\frac{(-1)^r\gamma_r}{r!}z^r,
 \qquad \gamma_0=\gamma.
\end{equation}
Fix $k\in\mathbb Z\setminus\{0\}$ and put $\lambda_k=2\pi i k$. Define
\begin{equation}\label{eq:Kintro}
 K=\gamma_1+\frac{(\gamma-1)^2+\zeta(2)+1}{2},\qquad
 b=1-\gamma-\frac{4K}{3}.
\end{equation}
For large positive $X$, the two roots described above satisfy
\begin{equation}\label{eq:headline}
 t_\pm=e^{-X}\left\{1\pm\sqrt{\frac{-2K\lambda_k}{X}}
       +\frac{b\lambda_k}{X}+O_k(X^{-3/2})\right\}.
\end{equation}
Both roots use the branch of $\Log t$ near the positive real axis; the two
signs do not come from assigning different logarithm branches to the same
input. They are two distinct solutions of one locally specified analytic
equation. At exactly integer order the small positive inverse instead has
size $e^{-X}/X$. Thus even the leading inverse scale is not uniform in a
complex neighborhood of the order.

The conclusions of this paper concern genuine convergent multivariable
Taylor or square-root charts. They do not assert Borel summability,
resurgence, or Stokes constants. Local branch interchange at a fold and
monodromy around $t=0$ are explicitly distinguished.

\subsection{What is inherited and what is developed}
Jonqui\`ere's continuation formula is the starting identity
\cite[Eq.~25.12.12]{dlmf-poly}. The Gamma Taylor series and the zeta Laurent
convention are classical \cite{dlmf-gamma,dlmf-zeta}. The mechanism by which a
simple critical point gives a square root is classical, and is already
used extensively in the transseries collection. The contribution here is
not a new version of the implicit-function theorem. It is the explicit
choice of a pole-cancelled family, the two complementary order scalings,
the Gamma--zeta critical-locus formulas, the exceptional integer-order
classification, and the arithmetic coefficients of the resulting inverse
sectors.

The reading and search scope is recorded in Section~\ref{sec:audit} and
\path{provenance/SOURCE_AUDIT.md}. In particular, the supplied proofs do not
assume the correctness of unreviewed repository arrivals.

\section{A pole-cancelled polylogarithm family}
\label{sec:normal}
For $n\ge1$ define
\begin{equation}\label{eq:Fdef}
 \F_{n,\eps}(t)=(-1)^n n!\left\{
 \Li_{n+1+\eps}(e^{-t})-
 \sum_{j=0}^{n-1}\zeta(n+1+\eps-j)\frac{(-t)^j}{j!}\right\}.
\end{equation}
We work initially on $0<|t|<2\pi$, $|\arg t|<\pi$, with the branch continued
from positive real $t$. Write $L=-\Log t$, and set
\begin{align}
 g_n(z)&=\Gamma(1-z)\prod_{j=1}^n(1+z/j)^{-1},\label{eq:gdef}\\
 Z(z)&=\zeta(1+z)-z^{-1},\qquad Z(0)=\gamma,\label{eq:Zdef}\\
 R_n(t,z)&=n!\sum_{j\ge0}
    \frac{(-1)^{j+1}\zeta(z-j)}{(n+j+1)!}t^j.\label{eq:Rdef}
\end{align}
Let $H_n^{(r)}=\sum_{j=1}^n j^{-r}$ and $H_n=H_n^{(1)}$.
When $n=0$ is used in a ladder formula, set $H_0^{(r)}=0$,
$g_0(z)=\Gamma(1-z)$, and interpret empty sums and products in the
usual way. The same normal-form and jet formulas then apply to
$\F_{0,\eps}(t)=\Li_{1+\eps}(e^{-t})$.

\begin{theorem}[Exact holomorphic cancellation]\label{thm:normal}
For $|\eps|<1$ and the stated $t$-branch,
\begin{align}
 \F_{n,\eps}(t)&=t^n\B_n(\eps,L)+t^{n+1}R_n(t,\eps),\label{eq:normal}\\
 \B_n(\eps,L)&=\frac{1-g_n(\eps)e^{-\eps L}}{\eps}+Z(\eps).
 \label{eq:block}
\end{align}
The quotient in \eqref{eq:block} is removable at $\eps=0$, and
\begin{equation}\label{eq:blockzero}
 \B_n(0,L)=L+H_n,\qquad R_n(0,\eps)=-\frac{\zeta(\eps)}{n+1},\qquad
 R_n(0,0)=\frac1{2(n+1)}.
\end{equation}
The series for $R_n$ converges normally on compact subsets of
$\{|t|<2\pi,\ |\eps|<1\}$. All fixed order derivatives in $\eps$ may be
taken term by term there.
\end{theorem}
\begin{proof}
For noninteger $s$, the classical continuation formula gives
\[
 \Li_s(e^{-t})=\Gamma(1-s)t^{s-1}
       +\sum_{j\ge0}\zeta(s-j)\frac{(-t)^j}{j!}.
\]
At $s=n+1+\eps$, the Gamma functional equation yields
\[
 (-1)^n n!\,\Gamma(-n-\eps)=-\frac{g_n(\eps)}{\eps}.
\]
After subtracting the first $n$ Taylor terms, the term with $j=n$ is
$t^n\zeta(1+\eps)$. Combining it with the Gamma term gives
\eqref{eq:block}; reindexing the remaining series gives \eqref{eq:Rdef}.
Since $g_n(0)=1$, its numerator vanishes at $\eps=0$. Also
$g_n'(0)=\gamma-H_n$, so its limiting value is
$L-(\gamma-H_n)+\gamma=L+H_n$.

For completeness, normal convergence does not follow merely from a
pointwise cancellation. On a compact set $|\eps|\le r<1$, the zeta
functional equation \cite{dlmf-functional} and the uniform Gamma-ratio
estimate give, for $j\ge2$,
\[
 \left|\frac{\zeta(\eps-j)}{(n+j+1)!}\right|
 \le C_r(2\pi)^{-j}j^{-n-1+r}.
\]
The two initial terms are holomorphic in the parameter disk. This bound
proves normal convergence when $|t|\le\rho<2\pi$. Cauchy's formula in a
slightly larger parameter disk then justifies all fixed parameter
derivatives. The equality at $\eps=0$ follows by removable continuation.
\end{proof}

\begin{proposition}[Differential ladder and positive integral]\label{prop:ladder}
With $\F_{0,\eps}(t)=\Li_{1+\eps}(e^{-t})$,
\begin{equation}\label{eq:ladder}
 \frac{\dd}{\dd t}\F_{n,\eps}(t)=n\F_{n-1,\eps}(t),\qquad
 \F_{n,\eps}(t)=n\int_0^t(t-u)^{n-1}\Li_{1+\eps}(e^{-u})\dd u.
\end{equation}
The integral holds for positive $t$ and $|\eps|<1$. For real $\eps$ in
this range, $\F_{n,\eps}$ is positive and strictly increasing on $t>0$.
Consequently every sufficiently small positive target has a unique
positive inverse.
\end{proposition}
\begin{proof}
Differentiate \eqref{eq:Fdef}, using
$\dd\Li_s(e^{-t})/\dd t=-\Li_{s-1}(e^{-t})$, and cancel the differentiated
Taylor polynomial. The constants at $t=0$ vanish for $n\ge1$ by
\eqref{eq:normal}, since $\Re\eps>-1$. Repeated integration gives the
second identity. Local integrability of the integrand follows from its
$u^{\eps}$ or logarithmic singularity. For real $\eps$, the defining
series of $\Li_{1+\eps}(e^{-u})$ consists of strictly positive terms.
The integral and the differential ladder give positivity and strict
monotonicity, including $n=1$.
\end{proof}

The subtraction in \eqref{eq:Fdef} is essential. For $n\ge2$, inversion of
the raw quantity $\zeta(n+1)-\Li_{n+1}(e^{-t})$ is ordinarily governed by a
nonzero linear term, not by $t^n\Log(1/t)$. This paper concerns a normalized
Taylor remainder, not an incorrectly assigned singular leading term for
the unsubtracted polylogarithm.

\section{All integer-order jets and their exact monodromy}
\label{sec:jets}
The Gamma expansion gives
\begin{equation}\label{eq:logg}
 \log g_n(z)=(\gamma-H_n)z+
 \sum_{j\ge2}\frac{\zeta(j)+(-1)^jH_n^{(j)}}j z^j,
 \qquad |z|<1.
\end{equation}
Define $a_n=\gamma-H_n$ and $b_{n,j}=\zeta(j)+(-1)^jH_n^{(j)}$.

\begin{theorem}[Complete logarithmic jet]\label{thm:jets}
For every $r\ge0$,
\begin{align}
 \left.\partial_\eps^r\F_{n,\eps}(t)\right|_{\eps=0}
 &=t^nP_{n,r}(L)+t^{n+1}R_{n,r}(t),\label{eq:jets}\\
 P_{n,r}(L)&=(-1)^r\gamma_r-r!\,[z^{r+1}]
 \exp\left((a_n-L)z+\sum_{j\ge2}\frac{b_{n,j}}jz^j\right),\label{eq:P}\\
 R_{n,r}(t)&=n!\sum_{j\ge0}
       \frac{(-1)^{j+1}\zeta^{(r)}(-j)}{(n+j+1)!}t^j.
 \label{eq:Rjet}
\end{align}
In particular $\deg_LP_{n,r}=r+1$, with leading coefficient
$(-1)^r/(r+1)$. Every coefficient in \eqref{eq:P} is a finite polynomial
in harmonic numbers, ordinary zeta values, Euler's constant, and one
Stieltjes constant.
\end{theorem}
\begin{proof}
Expand the removable quotient \eqref{eq:block}. Its coefficient of
$\eps^r$ is minus the coefficient of $z^{r+1}$ in $g_n(z)e^{-zL}$,
plus $(-1)^r\gamma_r/r!$. Equation~\eqref{eq:logg} yields \eqref{eq:P};
normal convergence proves \eqref{eq:Rjet}. The highest power of $L$ can
come only from $e^{-zL}$ and is
$-r!(-L)^{r+1}/(r+1)!=(-1)^rL^{r+1}/(r+1)$.
\end{proof}
The first three polynomials are
\begin{align}
 P_{n,0}(L)&=L+H_n,\label{eq:jetexamples0}\\
 P_{n,1}(L)&=-\gamma_1-\frac{(L-a_n)^2+b_{n,2}}2,\label{eq:jetexamples1}\\
 P_{n,2}(L)&=\gamma_2+\frac{(L-a_n)^3}{3}
          +b_{n,2}(L-a_n)-\frac{2b_{n,3}}3.\label{eq:jetexamples2}
\end{align}
The ladder also gives the exact finite-polynomial relation
\begin{equation}\label{eq:Pladder}
 nP_{n,r}(L)-P_{n,r}'(L)=nP_{n-1,r}(L).
\end{equation}
This can be proved either by differentiation of \eqref{eq:jets}, or from
$g_{n-1}(z)=(1+z/n)g_n(z)$. Its analytic tails satisfy the corresponding
coefficient identity.

\begin{proposition}[Sharp unipotent monodromy of every jet]\label{prop:monodromy}
Let $M$ be counterclockwise continuation once around $t=0$ within the
logarithmic covering of $0<|t|<2\pi$, and set $\Delta=M-1$. Then
\begin{equation}\label{eq:monodromyfull}
 (M-1)\F_{n,\eps}(t)
 =-t^{n+\eps}g_n(\eps)\frac{e^{2\pi i\eps}-1}{\eps}.
\end{equation}
If $J_{n,r}(t)=\left.\partial_\eps^r\F_{n,\eps}(t)\right|_0$, then
\begin{equation}\label{eq:monodromyjet}
 \Delta^{r+2}J_{n,r}=0,\qquad
 \Delta^{r+1}J_{n,r}=-r!(2\pi i)^{r+1}t^n\ne0.
\end{equation}
\end{proposition}
\begin{proof}
Only $t^{\eps}$ changes in \eqref{eq:normal}; this proves
\eqref{eq:monodromyfull}. For the jet, continuation sends
$L$ to $L-2\pi i$ and leaves the analytic tail unchanged. A finite
difference lowers the degree of a polynomial by one. The $(r+1)$st
difference of its leading term is
$(-1)^rr!(-2\pi i)^{r+1}=-r!(2\pi i)^{r+1}$.
\end{proof}
This is a direct consequence of classical continuation, not a claim of a
new polylogarithm monodromy theorem. Its useful feature here is that the
nilpotence index and top finite difference are exact, independent of $n$,
and compatible with every order derivative and normalized remainder.

\section{The primary, uniform inverse chart}
\label{sec:primary}
Let
\begin{equation}\label{eq:A}
 \A(\lambda)=\frac{1-e^{-\lambda}}{\lambda},\quad\A(0)=1,
 \qquad \lambda_k=2\pi i k\quad(k\in\mathbb Z\setminus\{0\}).
\end{equation}
The zeros of $\A$ are exactly the $\lambda_k$, and each is simple. On the
real line $\A(\lambda)=\int_0^1 e^{-u\lambda}\dd u>0$.
For a positive target $y$ put
\begin{equation}\label{eq:primaryscale}
 X=\frac1n\log(1/y),\quad v=X^{-1},\quad h=X^{-1}\log X,
 \quad\eps=\lambda/X,\quad
 \tau=e^{-X}\{X\A(\lambda)\}^{-1/n}.
\end{equation}
For complex $\lambda$ choose $\log\A$ on a simply connected domain avoiding
its zeros; the root in \eqref{eq:primaryscale} uses that choice. The
logarithmic sheet is specified by
$\Log t=-X-(\log X+\log\A(\lambda))/n+\log U$.
The theorem below is on this sheet; it is on the initial branch whenever
its imaginary part lies in $(-\pi,\pi)$. If the parameter domain contains zero, normalize $\log\A(0)=0$.
On real parameter intervals we take the real logarithm of the positive
function $\A$.

\begin{theorem}[Convergent order-crossover inversion]\label{thm:primary}
Fix $n\ge1$ and a compact subset $\mathscr K$ of such a parameter domain.
There is a holomorphic function $U_n(v,h,\tau;\lambda)$ in a common
polydisk around $(0,0,0)$, locally holomorphic in $\lambda$, such that
$U_n(0,0,0;\lambda)=1$ and
\begin{equation}\label{eq:primaryinverse}
 t=\tau U_n(v,h,\tau;\lambda)
\end{equation}
solves $\F_{n,\lambda/X}(t)=y$ after the substitutions
\eqref{eq:primaryscale}. It is the unique solution with $t/\tau$ sufficiently
close to $1$, uniformly for $\lambda\in\mathscr K$.
Its Taylor series is absolutely convergent after substitution and can be
grouped as
\begin{equation}\label{eq:primarysectors}
 U_n=\sum_{c\ge0}\tau^c\sum_{N\ge0}X^{-N}
           P_{N,c}(\log X;\lambda),\qquad \deg P_{N,c}\le N.
\end{equation}
The coefficients are holomorphic in $\lambda$, and $P_{0,c}=0$ for $c\ge1$.
For real $\lambda$ in a compact interval this is the unique small positive
inverse.
\end{theorem}
\begin{proof}
With $t=\tau U$ we have
\[
 L=X+\frac{\log X+\log\A(\lambda)}n-\log U.
\]
Define the holomorphic function
\begin{align}
 \Phi_n(U;v,h,\lambda)
 &=\frac{1-e^{-\lambda}g_n(\lambda v)
       \exp\{-\lambda h/n-\lambda v\log\A(\lambda)/n
                          +\lambda v\log U\}}{\lambda}
       +vZ(\lambda v),\label{eq:Phi}\\
 \E_n(U;v,h,\tau,\lambda)
 &=U^n\{\Phi_n(U;v,h,\lambda)
               +v\tau U R_n(\tau U,\lambda v)\}-\A(\lambda).
 \label{eq:E}
\end{align}
At $\lambda=0$ the quotient is removable as a function of \emph{all}
variables: its numerator vanishes identically when $\lambda=0$. The
limiting formula, which is sometimes computationally preferable, is
\begin{equation}\label{eq:Phi0}
 \Phi_n(U;v,h,0)=1+h/n+v(H_n-\log U).
\end{equation}
At $(v,h,\tau)=(0,0,0)$,
\[
 \E_n(U;0,0,0,\lambda)=\A(\lambda)(U^n-1),\qquad
 \partial_U\E_n(1;0,0,0,\lambda)=n\A(\lambda)\ne0.
\]
The analytic implicit-function theorem gives the stated root. Uniqueness
glues the local parameter charts. Compactness supplies common smaller
radii and a common bound. Substituting the physical values into
\eqref{eq:E} is exactly \eqref{eq:normal} divided by $y$; hence this is an
actual inverse, not only a formal solution.

Write its Taylor series as $\sum u_{abc}(\lambda)v^ah^b\tau^c$.
After substitution its individual monomials become
\[u_{abc}X^{-a-b}(\log X)^b\tau^c.\]
This proves the degree assertion and absolute regroupability. When $v=0$, equation \eqref{eq:E} is independent
of $\tau$, so $U_n(0,h,\tau;\lambda)$ is as well. In particular
$u_{00c}=0$ for $c>0$. For real parameters, conjugation and uniqueness make
the root real and close to one; positivity and the uniqueness assertion in
Proposition~\ref{prop:ladder} identify it with the positive inverse.
\end{proof}

\begin{corollary}[First crossover coefficient]\label{cor:primaryfirst}
Uniformly on the parameter compacts of Theorem~\ref{thm:primary},
\begin{equation}\label{eq:primaryfirst}
 \frac{t}{\tau}=1-
 \frac{e^{-\lambda}\{(\log X+\log\A(\lambda))/n-a_n\}+\gamma}
      {n\A(\lambda)X}
 +O\left(\frac{(\log X)^2}{X^2}+\frac{|\tau|}{X}\right).
\end{equation}
\end{corollary}
\begin{proof}
The terms of first degree in $v,h$ in \eqref{eq:Phi} are
$e^{-\lambda}h/n+
 v\{e^{-\lambda}(\log\A/n-a_n)+\gamma\}$.
Linearizing $U^n$ at one gives \eqref{eq:primaryfirst}. Holomorphy bounds
the remaining slow terms quadratically. The difference between the full
root and its value at $\tau=0$ is divisible by $v\tau$ and is therefore
$O(|\tau|/X)$ after substitution.
\end{proof}

\subsection{Remainders and exponential resolution}
If $U_n$ is bounded by $M$ on a closed polydisk of equal radius $\rho$,
Cauchy's estimate gives an entirely explicit form of the truncation control.
For $s=|v|+|h|+|\tau|<\rho$, truncation by total degree at most $N$ obeys
\begin{equation}\label{eq:cauchyrem}
 |U_n-U_{n,\le N}|\le M\frac{(s/\rho)^{N+1}}{1-s/\rho}.
\end{equation}
Indeed, the sum of absolute monomials of degree $d$ is at most $s^d$.
The radii are existential here; no interval-certified numerical radius is
claimed. Normal convergence is stronger than a bare Poincar\'e expansion.

A finite slow expansion has an algebraic error much larger than the first
exponential sector. Therefore a claim about that sector must refer to the
analytic coefficient of $\tau$, or to subtraction of the entire convergent
slow block, not to subtraction of a few powers of $1/X$. This distinction
will also be crucial at the fold.

\subsection{Comparison with the exact integer Lambert core}
At $\eps=0$, ignoring only the analytic tail, the equation is
$y=t^n(H_n-\log t)$. Its small positive core inverse is
\begin{equation}\label{eq:Lambertcore}
 t_{\rm core}=\left\{
 \frac{-ny}{W_{-1}(-ny e^{-nH_n})}\right\}^{1/n}.
\end{equation}
To verify this, put $w=n(H_n-\log t)$. Then
$we^{-w}=ny e^{-nH_n}$, so $-w=W_{-1}(-ny e^{-nH_n})$.
The negative branch gives $w\to+\infty$ and hence $t\to0$.
This is a classical Lambert reduction \cite{dlmf-w}, not a new special
function identity. The primary chart recovers its logarithmic expansion,
while retaining both order variation and the full analytic tail.

\section{Secondary resonances and an exact Gamma--zeta critical locus}
\label{sec:secondary}
Fix $\lambda_k=2\pi i k\ne0$. Near this value it is no longer appropriate
to divide the inverse scale by $\A(\lambda)$. Instead put
\begin{equation}\label{eq:secondaryscale}
 y=e^{-nX},\qquad \eta=e^{-X},\qquad t=\eta U,\qquad
 \lambda=\eps X,\qquad \kappa=\frac{\lambda-\lambda_k}{\eps}.
\end{equation}
The last identity is used when $\eps\ne0$; the resulting local equation
extends holomorphically to $\eps=0$. Since $e^{-\lambda}=e^{-\eps\kappa}$,
\eqref{eq:normal} becomes exactly
\begin{align}
 \HH_n(U;\eps,\kappa,\eta)
 &=U^n\left\{\frac{1-g_n(\eps)e^{\eps(\log U-\kappa)}}{\eps}
                     +Z(\eps)\right\}
       +\eta U^{n+1}R_n(\eta U,\eps)-1=0.\label{eq:H}\\
 \HH_n(U;0,\kappa,\eta)
 &=U^n(H_n+\kappa-\log U)
       +\eta U^{n+1}R_n(\eta U,0)-1.\label{eq:H0all}
\end{align}
Here $\log U$ is the local logarithm near a specified positive value of
$U$. The variables $(\eps,\kappa,\eta)$ are first independent analytic
coordinates. The original problem is recovered by \eqref{eq:secondaryscale}.
Keeping this distinction is essential: $\eta=0$ defines the analytic
\emph{singular core}, not a literal substitution into $e^{-X}$ at finite $X$.

\subsection{Critical points of the singular core at every order}
\begin{theorem}[Exact core critical locus]\label{thm:corecritical}
For each integer $n\ge1$, the singular core
$\HH_n(U;\eps,\kappa,0)$ has a unique nearby critical solution
$(U,\kappa)=(u_n(\eps),k_n(\eps))$ at which both the value and its $U$
derivative vanish, where
\begin{align}
 u_n(\eps)&=\left\{\frac{n+\eps}{\eps\zeta(1+\eps)}\right\}^{1/n},
       \label{eq:ucore}\\
 k_n(\eps)&=\frac{\log g_n(\eps)-\log n
                        +(n+\eps)\log u_n(\eps)}{\eps}.
       \label{eq:kcore}
\end{align}
The root and logarithms are chosen analytically near
$u_n(0)=n^{1/n}$. All displayed singularities at $\eps=0$ are removable,
and
\begin{align}
 k_n(0)&=\kappa_n^*:=\frac{1+\log n}{n}-H_n,\label{eq:kstar}\\
 k_n(\eps)&=\kappa_n^*+K_n\eps+O(\eps^2),\label{eq:kexpand}\\
 K_n&=\gamma_1+\frac{(\gamma-1/n)^2}{2}
                     +\frac{\zeta(2)+H_n^{(2)}}2>0.\label{eq:Kn}
\end{align}
At the core critical point, the following exact derivative identities hold:
\begin{equation}\label{eq:corederivatives}
 \partial_\kappa\HH_n=n,\qquad
 \partial_U^2\HH_n=-\frac{n(n+\eps)}{u_n^2},\qquad
 \partial_U\partial_\kappa\HH_n=\frac{n(n+\eps)}{u_n}.
\end{equation}
\end{theorem}
\begin{proof}
For $\eps\ne0$, let $Q=g_n(\eps)e^{-\eps\kappa}$. The core equation
and its derivative are
\[
 U^n\zeta(1+\eps)-\frac{Q}{\eps}U^{n+\eps}=1,\qquad
 nU^{n-1}\zeta(1+\eps)-\frac{n+\eps}{\eps}QU^{n+\eps-1}=0.
\]
Eliminating $Q$ gives $U^n=(n+\eps)/(\eps\zeta(1+\eps))$.
The same equations give $QU^{n+\eps}=n$, whose logarithm gives
\eqref{eq:kcore}. Conversely these two identities imply both equations,
so there is no lost condition. The analytic function
$\eps\zeta(1+\eps)$ equals $1$ at zero. This fixes the analytic root
and logarithms in a sufficiently small disk. The numerator of
\eqref{eq:kcore} vanishes at zero, which proves removability.

To obtain its first coefficients, use \eqref{eq:logg} and
\[
 \log u_n(\eps)=\frac{\log n}{n}
 +\left(\frac1{n^2}-\frac\gamma n\right)\eps
 +\left(-\frac1{2n^3}+\frac{\gamma_1}{n}
                            +\frac{\gamma^2}{2n}\right)\eps^2
 +O(\eps^3).
\]
Substitution gives \eqref{eq:kstar}--\eqref{eq:Kn}.
Direct differentiation, followed by
$u_n^n\zeta(1+\eps)=(n+\eps)/\eps$ and $Qu_n^{n+\eps}=n$,
gives \eqref{eq:corederivatives}. The Jacobian of
$(\HH_n,\partial_U\HH_n)$ with respect to $(\kappa,U)$ has determinant
$-n^2(n+\eps)/u_n^2\ne0$. The implicit-function theorem therefore proves
local uniqueness, including at $\eps=0$.

We give a nonnumerical proof of the positivity in \eqref{eq:Kn}.
Euler partial summation gives, first for $\Re s>1$,
\[
 \zeta(s)=\frac{s}{s-1}-s\int_1^\infty\{x\}x^{-s-1}\dd x.
\]
The integral is holomorphic for $\Re s>0$, and comparison at $s=1$ gives
\[
 \gamma_1=\int_1^\infty \{x\}(1-\log x)x^{-2}\dd x
 >-\int_{\mathrm e}^\infty(\log x-1)x^{-2}\dd x=-\mathrm e^{-1}.
\]
The strict inequality follows, for example, from the positive integral
on $(1,\mathrm e)$ and $0\le\{x\}<1$. Since $\zeta(2)>1$ and
$H_n^{(2)}\ge1$, we obtain $K_n>1-\mathrm e^{-1}>0$.
\end{proof}

\begin{theorem}[First deformation by the full analytic tail]
\label{thm:tailcritical}
There are unique holomorphic critical functions
$u_{n,c}(\eps,\eta)$ and $\kappa_{n,c}(\eps,\eta)$ near the core
locus, satisfying $\HH_n=\partial_U\HH_n=0$. Their expansions in $\eta$
are
\begin{align}
 \kappa_{n,c}(\eps,\eta)
 &=k_n(\eps)+\eta\frac{\zeta(\eps)u_n(\eps)^{n+1}}{n(n+1)}
                                      +O(\eta^2),\label{eq:kctail}\\
 u_{n,c}(\eps,\eta)
 &=u_n(\eps)-\eta\frac{(1-\eps)\zeta(\eps)u_n(\eps)^{n+2}}
                         {n(n+1)(n+\eps)}+O(\eta^2).
                                      \label{eq:uctail}
\end{align}
These remainders are uniform on smaller compact $\eps$-disks.
\end{theorem}
\begin{proof}
The nonzero Jacobian just computed applies unchanged at $\eta=0$, so the
holomorphic implicit-function theorem gives both functions. Write
$u_{n,c}=u_n+\eta u_1+O(\eta^2)$ and
$\kappa_{n,c}=k_n+\eta k_1+O(\eta^2)$, and put
$c=R_n(0,\eps)=-\zeta(\eps)/(n+1)$. The terms of first degree in
$\eta$ in the value and derivative equations are
\[
 nk_1+cu_n^{n+1}=0,
\qquad
 -\frac{n(n+\eps)}{u_n^2}u_1
 +\frac{n(n+\eps)}{u_n}k_1+(n+1)cu_n^n=0.
\]
Solving this triangular system gives \eqref{eq:kctail} and
\eqref{eq:uctail}. Holomorphy proves the uniform remainder assertion.
\end{proof}

\subsection{Why only the dilogarithm is an undetuned fold}
\begin{corollary}[Exceptional-order classification]\label{cor:exceptional}
Among $n\ge1$, the limiting secondary equation at $\kappa=0$ has a
positive double root precisely for $n=1$, when that root is $U=1$.
For every $n\ge2$ it has exactly two positive roots, both simple. Each
of those two simple roots extends to a holomorphic solution of the full
secondary equation near $(\eps,\kappa,\eta)=(0,0,0)$.
\end{corollary}
\begin{proof}
For positive $U$, the function $f(U)=U^n(H_n-\log U)$ increases up to
$U_*=e^{H_n-1/n}$ and decreases thereafter. It vanishes at the limiting
left endpoint and at $U=e^{H_n}$, and its maximum is
$e^{nH_n-1}/n$. This maximum equals $1$ at $n=1$.
For $n\ge2$, $nH_n>1+\log n$, so the maximum is strictly greater than
one. To see the strict inequality directly, $H_n\ge1$ and
$n>1+\log n$ for $n>1$ already suffice. This proves the root count and
simplicity. The last claim is the implicit-function theorem at each root.
Equivalently, the critical detuning \eqref{eq:kstar} is zero only at
$n=1$, and is negative for all larger $n$.
\end{proof}

The limiting roots also have a useful classical Lambert description. For
real $\kappa>\kappa_n^*$ they are
\begin{equation}\label{eq:secondaryLambert}
 U_j(\kappa)=\left\{
 \frac{-n}{W_j(-n e^{-n(H_n+\kappa)})}\right\}^{1/n},
 \qquad j\in\{0,-1\}.
\end{equation}
Indeed, substituting $w=n(H_n+\kappa-\log U)$ reduces the equation to
$we^{-w}=ne^{-n(H_n+\kappa)}$. At $\kappa=\kappa_n^*$ the Lambert
argument equals $-1/\mathrm e$ and the two branches meet. As real
$\kappa\to+\infty$, the $W_{-1}$ root satisfies $U\sim\kappa^{-1/n}$,
whereas the $W_0$ root satisfies $U\sim e^{H_n+\kappa}$.
These are identities for the limiting secondary core, not a proof of a
global connection between every finite-$X$ primary and secondary branch.

\section{The dilogarithm: exact critical curve and a two-sheet chart}
\label{sec:fold}
From now on $n=1$, unless indicated otherwise. Set
\begin{equation}\label{eq:abK}
 a=\gamma-1,\qquad
 K=\gamma_1+\frac{a^2+\zeta(2)+1}{2},\qquad
 b=-a-\frac{4K}{3}.
\end{equation}
We suppress $n=1$ in $\HH$, $u_c$ and $\kappa_c$.

\subsection{A closed Gamma--zeta expression}
Theorem~\ref{thm:corecritical} specializes to
\begin{align}
 u_1(\eps)&=\frac{1+\eps}{\eps\zeta(1+\eps)},\label{eq:udilog}\\
 k_1(\eps)&=\log(1+\eps)+
 \frac{\log\Gamma(1-\eps)
                  -(1+\eps)\log(\eps\zeta(1+\eps))}{\eps}.
                                                        \label{eq:kdilog}
\end{align}
The logarithms in \eqref{eq:kdilog} vanish at zero. Expanding gives
\begin{align}
 u_1(\eps)&=1-a\eps+O(\eps^2),\label{eq:udilogfirst}\\
 k_1(\eps)&=K\eps+J\eps^2+O(\eps^3),\label{eq:kdilogsecond}\\
 J&=-\frac12+\frac{\zeta(3)}3+(1-\gamma)\gamma_1
          +\frac{\gamma^2}{2}-\frac{\gamma^3}{3}-\frac{\gamma_2}{2}.
                                                        \label{eq:J}
\end{align}
For example,
$\log(\eps\zeta(1+\eps))=
 \gamma\eps-(\gamma_1+\gamma^2/2)\eps^2+
 (\gamma_2/2+\gamma\gamma_1+\gamma^3/3)\eps^3+O(\eps^4)$,
which verifies every sign in \eqref{eq:J} directly.
The full critical curve satisfies
\begin{equation}\label{eq:dilogcriticalfirst}
 \kappa_c(\eps,\eta)=K\eps-\eta/4+\Ot(2),\qquad
 u_c(\eps,\eta)=1-a\eps+\eta/4+\Ot(2).
\end{equation}
Here $\Ot(d)$ denotes a convergent analytic series whose monomials have
total degree at least $d$ in the displayed small parameters.

\subsection{Exact roots from analytic root moments}
\begin{theorem}[Convergent two-sheet normal form]\label{thm:twosheet}
There are holomorphic functions $C(\eps,\kappa,\eta)$ and
$D(\eps,\kappa,\eta)$ near zero such that the full equation
$\HH(U;\eps,\kappa,\eta)=0$ has exactly two roots in a fixed sufficiently
small disk about $U=1$, counted with multiplicity, and they are exactly
\begin{equation}\label{eq:CD}
 U_\pm=C(\eps,\kappa,\eta)\pm\sqrt{D(\eps,\kappa,\eta)}.
\end{equation}
Their first analytic jets are
\begin{align}
 C&=1+\frac43\kappa+b\eps+\frac7{12}\eta+\Ot(2),\label{eq:Cjet}\\
 D&=2\kappa-2K\eps+\eta/2+\Ot(2).\label{eq:Djet}
\end{align}
Moreover
\begin{equation}\label{eq:Dfactor}
 D=(\kappa-\kappa_c(\eps,\eta))d(\eps,\kappa,\eta),
 \qquad d(0,0,0)=2,
\end{equation}
where $d$ is holomorphic and nonvanishing after shrinking the neighborhood.
Thus a small loop in $\kappa$ around its critical value, with the other
parameters fixed, interchanges the two roots.
\end{theorem}
\begin{proof}
At zero parameters,
\[
 H_0(U):=\HH(U;0,0,0)=U(1-\log U)-1
       =-\frac{(U-1)^2}{2}+O((U-1)^3).
\]
Choose a circle $\Gamma$ about $1$ so small that this is its only zero in
the closed disk. The function is nonzero on $\Gamma$. Holomorphy and
Rouch\'e's theorem show that the full equation has exactly two zeros
inside, counted with multiplicity, for all sufficiently small parameters.
Define
\[
 S_j=\frac1{2\pi i}\int_\Gamma U^j
               \frac{\partial_U\HH}{\HH}\dd U\quad(j=1,2),\qquad
 C=S_1/2,\qquad D=S_2/2-C^2.
\]
The contour is fixed and its denominator stays nonzero, so the moments
are holomorphic even when the roots coalesce. If the two roots are
$u_+,u_-$, the argument principle gives
$S_j=u_+^j+u_-^j$. Hence $C=(u_++u_-)/2$ and
$D=(u_+-u_-)^2/4$, proving \eqref{eq:CD} without a prior choice of labels.

For the first jets put $\delta=U-1$. Then
\[
 \frac1{H_0}=-2\delta^{-2}-\frac23\delta^{-1}+O(1).
\]
The first parameter-degree part $\Delta_1$ of $\HH-H_0$ is
\[
 \Delta_1=\kappa U-\eps U
       \{K+a\log U+(\log U)^2/2\}+\eta U^2/4.
\]
Write $\Delta_1=q_0+q_1\delta+O(\delta^2)$, where
$q_0=\kappa-K\eps+\eta/4$ and
$q_1=\kappa-(K+a)\eps+\eta/2$.
Expanding the logarithmic derivative and integrating by parts shows that
the first variations are
\[
 S_1-2=2q_1+\frac23q_0,\qquad
 S_2-2=4q_1+\frac{16}3q_0.
\]
Consequently $C-1=q_1+q_0/3$ and $D=2q_0$ to first degree. These are
\eqref{eq:Cjet}--\eqref{eq:Djet}. In particular $\partial_\kappa D(0)=2$.
A zero of $D$ is precisely a double root, and its locus is the unique
critical curve in Theorem~\ref{thm:tailcritical}. Division by
$\kappa-\kappa_c$ therefore gives \eqref{eq:Dfactor}; for example the
quotient is the integral of $\partial_\kappa D$ along the segment from
$\kappa_c$ to $\kappa$. The nonzero analytic factor has a local analytic
square root. The remaining square root changes sign on one loop, which
proves branch interchange.
\end{proof}

This local interchange is ordinary ramification of an analytic inverse.
It is not a Stokes jump. Separately, Proposition~\ref{prop:monodromy}
concerns continuation of the forward family around $t=0$. Conflating these
two operations would give incorrect branch bookkeeping.

\subsection{A moving fold in the actual order parameter}
\begin{theorem}[Moving physical critical order]\label{thm:movingfold}
Fix $k\ne0$ and write $v=X^{-1}$. There is a holomorphic function
$\lambda_c(v,\eta)$ near $(0,0)$ with $\lambda_c(0,0)=\lambda_k$,
characterized by
\begin{equation}\label{eq:lambdacimplicit}
 \lambda_c-\lambda_k
       =\lambda_cv\,\kappa_c(\lambda_cv,\eta).
\end{equation}
Its expansion is
\begin{equation}\label{eq:lambdac}
 \lambda_c=\lambda_k+K\lambda_k^2v^2
       -\frac{\lambda_k}{4}v\eta
       +O(v^3+v^2\eta+v\eta^2).
\end{equation}
On the physical specialization $\eta=e^{-X}$, set
$\eps_c=\lambda_c/X$ and
$t_c=\eta u_c(\eps_c,\eta)$. Then
\begin{equation}\label{eq:physicalfold}
 \zeta(2+\eps_c)-\Li_{2+\eps_c}(e^{-t_c})=e^{-X},
 \qquad \Li_{1+\eps_c}(e^{-t_c})=0,
\end{equation}
and this is a nondegenerate critical point of the inverse problem. In
particular,
\begin{align}
 \eps_c&=\frac{\lambda_k}{X}+\frac{K\lambda_k^2}{X^3}
                      +O_k(X^{-4}+e^{-X}X^{-2}),\label{eq:epsc}\\
 t_c&=e^{-X}\{1+(1-\gamma)\lambda_k/X
                               +O_k(X^{-2}+e^{-X})\}.\label{eq:tc}
\end{align}
\end{theorem}
\begin{proof}
The left side of \eqref{eq:lambdacimplicit} minus its right side is
holomorphic in $(\lambda,v,\eta)$, vanishes at $(\lambda_k,0,0)$, and
has $\lambda$ derivative one there. The implicit-function theorem proves
existence and uniqueness. Insert \eqref{eq:dilogcriticalfirst} and use
$\lambda_c=\lambda_k+O(v^2+v\eta)$ to obtain \eqref{eq:lambdac}.
For physical parameters, \eqref{eq:secondaryscale} makes the critical
secondary equation identical to the original normalized equation. Its
$U$ derivative is the $t$ derivative of $\F_{1,\eps}(t)$, which equals
$\Li_{1+\eps}(e^{-t})$. Its second $U$ derivative remains close to $-1$
and is nonzero. This proves \eqref{eq:physicalfold} and nondegeneracy.
Finally substitute \eqref{eq:lambdac} into \eqref{eq:dilogcriticalfirst}
and divide by $X$ where appropriate.
\end{proof}

The vanishing $X^{-1}$ displacement in \eqref{eq:lambdac} is another
expression of the exceptional-order classification. For general $n$, the
same argument at the tuned critical detuning gives instead
$\lambda_{n,c}=\lambda_k+\lambda_k\kappa_n^*v+O(v^2+v\eta)$.
This follows directly from \eqref{eq:kstar}; no analogous cancellation
occurs for $n\ge2$.

\section{Convergent secondary transseries and every exponential sector}
\label{sec:sectors}
Fix the undetuned resonance $\lambda=\lambda_k$ and hence $\kappa=0$.
The small variables that preserve the fold splitting are
\begin{equation}\label{eq:wr}
 v=X^{-1},\qquad w=\sqrt v,\qquad
 \eta=e^{-X},\qquad r=\eta/v=Xe^{-X},\qquad A_k=-2K\lambda_k\ne0.
\end{equation}
The letter $A_k$ here is a nonzero constant, not the function $\A$.

\begin{theorem}[Two-generator convergent inverse transseries]
\label{thm:secondarytransseries}
There are holomorphic functions $C_k(v,r)$ and $B_k(v,r)$ in a polydisk
about zero such that, for all sufficiently large real $X$, the two roots
near $t=e^{-X}$ are exactly
\begin{equation}\label{eq:secondarytransseries}
 t_\pm=e^{-X}\{C_k(X^{-1},Xe^{-X})
               \pm X^{-1/2}B_k(X^{-1},Xe^{-X})\}.
\end{equation}
One may select either square root of $A_k$, after which
\begin{equation}\label{eq:CBboundary}
 C_k(0,r)=1,\qquad B_k(0,r)=\sqrt{A_k+r/2}.
\end{equation}
Consequently \eqref{eq:secondarytransseries} is an absolutely convergent
transseries on the finite positive monomial grid generated by
$X^{-1/2}$ and $Xe^{-X}$, with the prefactor $e^{-X}$. Both roots use the
same original logarithm branch, and they are distinct for large $X$.
\end{theorem}
\begin{proof}
Substitute $(\eps,\kappa,\eta)=(\lambda_kv,0,vr)$ into
Theorem~\ref{thm:twosheet}. Define
\[
 C_k(v,r)=C(\lambda_kv,0,vr),\qquad
 T_k(v,r)=\frac{D(\lambda_kv,0,vr)}v.
\]
Every monomial in the numerator of the latter expression is divisible by
$v$, because $D(0,0,0)=0$. Thus $T_k$ is holomorphic, not just bounded
on a punctured set. Equation~\eqref{eq:Djet} gives
$T_k(0,r)=A_k+r/2$. Since $A_k\ne0$, shrink the polydisk and take its
analytic square root $B_k$. This proves \eqref{eq:CBboundary} and the
exact formula. The functions $C_k(w^2,r)\pm wB_k(w^2,r)$ are holomorphic
in $(w,r)$, so their Taylor series converge absolutely on a smaller
polydisk. Both generators tend to zero after \eqref{eq:wr}; this proves
convergence on the physical path. For each exponential action the powers
of $X^{-1/2}$ have a lower bound, giving a well-based monomial support.
Since $T_k\to A_k\ne0$, the roots stay distinct. Finally $U_\pm\to1$,
so $t_\pm$ remain in a narrow sector about the positive real axis and
share the logarithm specified at the start of the paper.
\end{proof}

A two-variable version of \eqref{eq:cauchyrem} gives a geometric error
bound in $|w|+|r|$. More precisely, if the analytic root expression is
bounded by $M$ on a closed polydisk of radius $\rho$, then its
Taylor polynomial of total degree at most $N$ has error at most
$M\theta^{N+1}/(1-\theta)$, where
$\theta=(|w|+|r|)/\rho<1$. Multiplication by $e^{-X}$ gives the error for
$t$. The statement is local; the paper does not supply a maximal numerical
polydisk or claim uniformity when $k$ grows with $X$.

\subsection{The first three algebraic corrections}
Put
\begin{align}
 J_0&=\frac{\gamma_2}{2}-\frac{a^3}{6}
            -\frac{a(\zeta(2)+1)}2-\frac{\zeta(3)-1}{3},\label{eq:J0}\\
 M_0&=\frac{a^2}{2}+(3a+1)K+\frac{13K^2}{9}+J_0.\label{eq:M0}
\end{align}
(The subscript distinguishes $M_0$ from a Cauchy bound.)
\begin{corollary}[Explicit Puiseux coefficients]\label{cor:puiseux}
At fixed $k\ne0$ and with a fixed choice of $\sqrt{A_k}$,
\begin{equation}\label{eq:threecorrections}
 \frac{t_\pm}{e^{-X}}=
 1\pm\frac{\sqrt{A_k}}{X^{1/2}}+\frac{b\lambda_k}{X}
 \pm\frac{M_0\lambda_k^2}{\sqrt{A_k}\,X^{3/2}}
 +O_k(X^{-2}).
\end{equation}
\end{corollary}
\begin{proof}
First set $\eta=0$ and write $\delta=U-1$. Give $\delta$ weight one
and $\eps$ weight two. Expansion of \eqref{eq:H} gives, through weight
four,
\begin{align*}
 \HH(1+\delta;\eps,0,0)
 ={}&-\delta^2/2+\delta^3/6-\delta^4/12\\
 &-\eps\{K+(K+a)\delta+(a+1)\delta^2/2\}
       +\eps^2J_0+O_{\mathrm{wt}}(5).
\end{align*}
Here the omitted analytic series has weight at least five. Substitute
$\eps=h^2$ and $\delta=ch+bh^2+dh^3+O(h^4)$.
The coefficients of $h^2,h^3,h^4$ give successively
$c^2=-2K$, $b=-a-4K/3$, and $cd=M_0$. This is the displayed formula
with $\eps=\lambda_k/X$. Convergence and the $O(h^4)$ error follow
from Theorem~\ref{thm:secondarytransseries} at $r=0$, rather than from
formal coefficient matching alone.
For the full equation, analyticity in $r$ shows that the additional
change in $U$ is $O(\sqrt v\,r+v r)$, hence
$O(e^{-X}\sqrt X)$ on the physical path. This is smaller than $X^{-2}$
and proves \eqref{eq:threecorrections} for the actual roots.
\end{proof}

\subsection{The leading coefficient of every exponential action}
The coefficients of the exponentially small sectors are most cleanly
specified before substitution. Write
\begin{equation}\label{eq:sectorfunctions}
 C_k(v,r)\pm\sqrt v\,B_k(v,r)
          =\sum_{m\ge0}r^m U_{\pm,m}(v).
\end{equation}
Each $U_{\pm,m}(w^2)$ is holomorphic in $w$. This independent-variable
definition removes the ambiguity inherent in trying to read an exponential
term beneath an unremoved algebraic error.

\begin{theorem}[All-action leading law]\label{thm:allsectors}
For every fixed integer $m\ge1$,
\begin{equation}\label{eq:Umleading}
 U_{\pm,m}(v)=
 \pm\binom{1/2}{m}2^{-m}A_k^{1/2-m}\sqrt v+O_{k,m}(v).
\end{equation}
Thus the sector of $t_\pm$ of exponential action $m+1$, as defined by
\eqref{eq:sectorfunctions}, has leading term
\begin{equation}\label{eq:allsectors}
 t_\pm^{[m]}=
 \pm\binom{1/2}{m}2^{-m}A_k^{1/2-m}
 e^{-(m+1)X}X^{m-1/2}\{1+O_{k,m}(X^{-1/2})\}.
\end{equation}
Every one of these coefficients is nonzero. In particular, the first two
additional sectors begin with
\begin{equation}\label{eq:firstactions}
 t_\pm^{[1]}\sim\pm\frac{e^{-2X}X^{1/2}}{4\sqrt{A_k}},\qquad
 t_\pm^{[2]}\sim\mp\frac{e^{-3X}X^{3/2}}{32A_k^{3/2}}.
\end{equation}
The full sector sum is absolutely convergent for all sufficiently large
$X$; \eqref{eq:allsectors} itself is an asymptotic statement at fixed $m$.
\end{theorem}
\begin{proof}
At $v=0$, $C_k$ is independent of $r$, so
$[r^m]C_k(v,r)=O(v)$ for $m\ge1$. Equation~\eqref{eq:CBboundary} and
Cauchy's formula in $r$ give
\[
 [r^m]B_k(v,r)=
 \binom{1/2}{m}2^{-m}A_k^{1/2-m}+O(v).
\]
This proves \eqref{eq:Umleading}. Now multiply by $e^{-X}r^m$ and
use $r=Xe^{-X}$ to obtain \eqref{eq:allsectors}. A half-integer binomial
coefficient $\binom{1/2}{m}$ never vanishes, and $A_k\ne0$; this proves
nonvanishing. The binomial coefficients at $m=1,2$ give
\eqref{eq:firstactions}. Absolute convergence is the two-variable Taylor
convergence already proved, with regrouping justified by absolute values.
\end{proof}

\begin{corollary}[Analytic-tail universality of the leading actions]
\label{cor:tailuniversality}
Replace $R_1(t,\eps)$ by any holomorphic function $\widetilde R(t,\eps)$
with $\widetilde R(0,0)=c$, while leaving the singular block unchanged.
All the local convergence statements remain valid. When $c\ne0$,
replace $2^{-m}$ by $(2c)^m$ in
\eqref{eq:Umleading}--\eqref{eq:allsectors}; all additional leading
sectors remain nonzero. When $c=0$, only the corresponding leading
coefficient in \eqref{eq:Umleading} vanishes: the relative asymptotic
formula \eqref{eq:allsectors} is not asserted, and a finer tail scaling
may be needed.
\end{corollary}
\begin{proof}
In the first-jet calculation the tail perturbation becomes
$\eta cU^2$, so $D_\eta(0)=2c$. Consequently
$T_k(0,r)=A_k+2cr$. The same analytic square root and coefficient argument
apply without any other change. This proves the claim, including its
qualification when $c=0$.
\end{proof}

\section{Finite residue formulas for every chart coefficient}
\label{sec:coefficients}
The convergent existence theorems admit a finite coefficient algorithm.
Although residue inversion is classical, the following formulas make the
present three-variable charts directly computable and avoid an unbounded
formal recursion for any fixed request.

\begin{theorem}[Finite coefficient extraction for the simple chart]
\label{thm:residuesimple}
Fix $n,\lambda$ in the setting of Theorem~\ref{thm:primary}. Write
$p=(v,h,\tau)$,
$E_0(U)=\A(\lambda)(U^n-1)$, and
$Q(U,p)=(\E_n(U,p,\lambda)-E_0(U))/E_0(U)$.
For every nonzero multi-index $\alpha$ of total degree $N$,
\begin{equation}\label{eq:residuesimple}
 [p^\alpha]U_n(p;\lambda)=
 \sum_{j=1}^{N}\frac{(-1)^j}{j}
       \Res_{U=1}[p^\alpha]Q(U,p)^j.
\end{equation}
Each residue requires only finitely many Taylor coefficients of the
explicit functions in \eqref{eq:E}.
\end{theorem}
\begin{proof}
Choose a small circle about $U=1$ containing no other zero of $E_0$.
For small $p$, $|Q|<1$ on the circle, and the simple root inside is
\[
 U_n(p)=\frac1{2\pi i}\int U\frac{\partial_U\E_n}{\E_n}\dd U
 =1-\frac1{2\pi i}\int\log(1+Q)\dd U.
\]
For the last equality, integrate the derivative of $U\log(1+Q)$ along
the closed circle; this logarithm is single valued there by its convergent
power series. Expanding the logarithm gives \eqref{eq:residuesimple}.
Since $Q$ has positive parameter degree, only $j\le N$ can contribute.
The pole order in the $j$th term is at most $j$, so a finite $U$-jet of
its analytic numerator determines the residue. All expansions and
integrals are uniformly convergent on the fixed contour for sufficiently
small $p$, justifying coefficient extraction.
\end{proof}

\begin{theorem}[Finite coefficient extraction at the double root]
\label{thm:residuedouble}
For the dilogarithm let $p=(\eps,\kappa,\eta)$,
$H_0(U)=U(1-\log U)-1$, and
$Q=(\HH-H_0)/H_0$. The root moments in
Theorem~\ref{thm:twosheet} satisfy, for $\alpha\ne0$ and
$N=|\alpha|$,
\begin{equation}\label{eq:residuedouble}
 [p^\alpha]S_\ell=
 \ell\sum_{j=1}^{N}\frac{(-1)^j}{j}
       \Res_{U=1}U^{\ell-1}[p^\alpha]Q(U,p)^j,
 \qquad \ell=1,2.
\end{equation}
Together with $C=S_1/2$ and $D=S_2/2-C^2$, this determines every
coefficient of the double-root chart from finite Taylor data. Substitution,
analytic square-root extraction, and ordinary convolution then determine
every coefficient in Theorem~\ref{thm:secondarytransseries}.
\end{theorem}
\begin{proof}
Use the same contour and logarithmic expansion as above, but start from
$S_\ell=(2\pi i)^{-1}\int U^\ell\HH_U/\HH\dd U$.
Its base value is $2$, because the base root at one has multiplicity two.
Integration by parts introduces $-\ell U^{\ell-1}$, giving the stated
formula. Only $j\le N$ contributes. The pole order is now at most $2j$,
which is still finite. Finally the analytic square root has nonzero
constant term $\sqrt{A_k}$, so its coefficient recurrence divides only by
$2\sqrt{A_k}$, never by a vanishing quantity.
\end{proof}

For the forward jets one can avoid even residues. If
$g_n(z)=\sum_{m\ge0}g_{n,m}z^m$ and
$\log g_n(z)=\sum_{j\ge1}\ell_{n,j}z^j$, then
\begin{equation}\label{eq:grecurrence}
 g_{n,0}=1,\qquad
 m g_{n,m}=\sum_{j=1}^{m}j\ell_{n,j}g_{n,m-j},
\end{equation}
where $\ell_{n,1}=a_n$ and $j\ell_{n,j}=b_{n,j}$ for $j\ge2$.
Multiplication by the finite jet of $e^{-Lz}$ gives
\eqref{eq:P}. This recurrence is obtained by differentiating $g_n=e^{\log
g_n}$ and comparing coefficients, so it is exact over a symbolic
coefficient ring. The verification program uses this recurrence rather
than numerical fitting of constants.

\section{Reproducible exact checks and numerical diagnostics}
\label{sec:verification}
The accompanying computations have deliberately different evidentiary
roles. The proofs above establish the analytic identities, convergence,
and root counts. Exact symbolic computations check a finite set of
coefficient identities. Floating-point evaluations test formulas and
branch selections at specified sample points; they do not certify the
last displayed digit, convergence radii, or publication novelty.

\subsection{Exact algebra}
The program \path{code/verify_exact.py}, using SymPy 1.14.0, completed
\textbf{150 exact assertions}. Its scopes are: the leading logarithmic
coefficient and exact monodromy degree for $0\le n,r\le4$; the differential
jet ladder; the Gamma recurrence through degree six; the first root-moment
coefficients of $C,D$; the Puiseux equation through fourth degree in
$\sqrt\eps$; the algebraic elimination and first-tail displacement at a
general core critical point; the constants $K,J$; and the all-sector
binomial convolution through degree eighteen.

Euler's constant, zeta values, and Stieltjes constants are independent
symbols in these polynomial checks. No conjectured numerical relation
between them is used. This exact algebra is neither a proof assistant
formalization nor an independent verification of complex-analytic
existence theorems. The complete receipt is
\path{results/exact_checks.json}; the symbolic output is
\path{results/symbolic_coefficients.json}.

\subsection{Floating-point comparisons}
The program \path{code/verify_numeric.py}, using mpmath 1.3.0 at
\textbf{180 decimal digits}, completed \textbf{73 diagnostics}.
It compares the pole-cancelled normal form with direct polylogarithm
calls, including integer and complex orders; checks first order jets by
a complex Cauchy coefficient calculation; solves primary inverses;
checks the closed core critical formulas and their full-tail deformation;
and checks both secondary roots and moving physical critical points
against direct polylogarithms. Several comparisons share the same numerical
library, so ``independent'' here means different evaluation routes, not
independent software implementations or interval arithmetic.

Direct subtraction of two $O(1)$ values to recover a quantity of size
$e^{-X}$ costs approximately $X/\log 10$ decimal digits. The program
therefore uses a large working precision and explicitly relaxes the
direct residual threshold in its most ill-conditioned sample, $X=320$.
It records every tolerance and measured residual in
\path{results/numeric_checks.json}. Finite tail truncations used in
these numerical comparisons are diagnostics, not interval bounds for the
omitted tail.

For the plus branch at $k=1$, let $U^{(2)}$ retain the square-root and
$X^{-1}$ corrections in \eqref{eq:threecorrections}, and let $U^{(3)}$
also retain the $X^{-3/2}$ correction. The observed errors are
\begin{center}
\begin{tabular}{rrrr}
\toprule
$X$ & $|U_+-U^{(2)}|$ & $|U_+-U^{(3)}|$ & $X^{3/2}|U_+-U^{(2)}|$\\
\midrule
80 & 0.0425674 & 0.00810337 & 30.4587 \\
160 & 0.0145921 & 0.00191344 & 29.5324 \\
320 & 0.00503756 & 0.000458763 & 28.8367 \\
\bottomrule
\end{tabular}

\end{center}
These values are rounded diagnostics, not rigorously enclosed decimals.
The constants used were approximately
\[
 K=1.3390244849427629699282971855433,\qquad
 b=-1.3625816448252168205109083374735.
\]
Their positivity or nonvanishing assertions in the theorems rely on exact
arguments, not these decimal values.

\subsection{How to replay}
Run the following commands from the package directory:
\begin{verbatim}
python code/verify_exact.py
python code/verify_numeric.py --dps 180
pdflatex -interaction=nonstopmode -halt-on-error article.tex
pdflatex -interaction=nonstopmode -halt-on-error article.tex
\end{verbatim}
The scripts need only the pinned Python dependencies in
\path{requirements.txt} and no network access after installation. The PDF
build uses the included numerical table; replaying the numerical script
regenerates it. The supplied \path{Makefile} separates mathematical replay
from PDF compilation. No repository mutation is performed by these files.

\section{Repository audit, corrections, and the novelty boundary}
\label{sec:audit}
This is a targeted research audit, not a review of every statement in
either large repository area. The observed \texttt{main} reference during
the session was
\begin{center}\small
\texttt{9ec0f0406afe7a5475b03ff65ed91f7e283979b8}.
\end{center}
Some initial index reads used the moving branch; later source and ledger
reads were explicitly pinned to that commit. The source manifest records
those distinctions and the returned blob identifiers rather than
pretending that all reads formed one atomic checkout.

\subsection{Overlap that must be credited}
The critical-Hahn article \cite{repo-hahn} already constructs an actual
convergent chart $x=T V(T,T^{2-\alpha})$ for a fixed exponent
$1<\alpha<2$, with finite positive support generators, coefficient
formulas, and a Cauchy remainder estimate. The inspected proof passage
uses the analytic implicit-function theorem correctly for that fixed
parameter regime. The present work does not supersede it and does not
identify a defect in it. Our varying-order resonance scale and the exact
Gamma--zeta critical locus are different statements. Uniformity through
an integer endpoint cannot be inferred from a fixed-parameter radius;
it requires the separate cancellation and scaling used here.

The transseries index also records existing moving-fold and inverse
branch-gluing work. Accordingly, a square-root chart, root moments,
Rouch\'e's theorem, finite residue inversion, and local branch interchange
are methods, not independent novelty claims of this paper. The specific
resonance lattice, critical-detuning classification, closed critical locus,
and the resulting sector coefficients are the candidate additions.

\subsection{A historical overstatement already corrected in the ledger}
The historical \emph{Polylogarithm--Polygamma Bridge} abstract
\cite{repo-bridge} uses stronger language about exact denominator
classification and independent new constants than finite numerical
searches establish. The source itself also distinguishes numerical
independence evidence from proof. Crucially, the canonical editorial
ledger \cite{repo-ledger} already directs replacement of the
``exactly crystallographic'' and independent-new-atom conclusions by
exact reductions and bounded searches. This is therefore \emph{not a
newly discovered uncorrected theorem error}.

The integration proposal is to preserve that existing correction and,
where standalone historical files are distributed, add an explicit
supersession pointer to the canonical ledger. No numerical
non-reduction should be promoted to a theorem of independence. This
paper neither needs nor supplies any such independence theorem.

\subsection{Preventive qualifications for the present contribution}
Several tempting but incorrect simplifications are ruled out by the
proofs. One must differentiate the pole-cancelled block, not two
individually singular terms evaluated at the integer. The primary inverse
must exclude the nonzero zeros of $\A$ or use a secondary chart.
An analytic-core critical identity must not be labeled exact for the full
polylogarithm without the tail deformation. Exponential coefficients must
be extracted from analytic sector functions, not fitted below a fixed
algebraic remainder. The two kinds of monodromy discussed above must
remain separate. Finally, the positivity and small-real-inverse statements
apply to the specified real family and do not exclude additional complex
roots outside the local disks.

No new false mathematical assertion in the inspected \emph{current
canonical} material was established in this audit. Its limited scope
would not justify claiming that the remaining corpus is error-free.
The paper supplies self-contained proofs for the new results so that
their validity does not depend on accepting unaudited sources.

\subsection{Publication-priority limitations}
A bounded web search for varying-order polylogarithm inversion,
Gamma--zeta resonance, Stieltjes/fold combinations, and polylogarithm
critical loci did not establish an exact antecedent for the combined
formulas here. Most relevant verified public background was classical
continuation and Lambert inversion in the DLMF. This limited search is
not a bibliographic proof of originality. In particular, local fold
methods are widespread, and equivalent formulas may occur in different
notation. Independent mathematical review and a broader literature
comparison are required before making a publication-priority claim.

\section{Further research questions and concrete next steps}
\label{sec:future}
\subsection{Matching the primary and secondary charts}
A local primary chart excludes $\lambda_k$, while the secondary chart
keeps $\kappa=(\lambda-\lambda_k)/\eps$ bounded. Construct an overlapping
atlas when $1\ll|\kappa|\ll X$, with a uniform remainder and an explicit
identification of which secondary root continues the small primary
inverse along a specified path. Formula~\eqref{eq:secondaryLambert}
suggests the matching law but does not prove the finite-$X$ analytic
continuation. The path dependence and possible encounters with the moving
fold should be part of the statement, not hidden in a branch convention.

\subsection{A growing resonance index}
All secondary theorems here fix $k$. Determine the largest regime
$k=k(X)$ for which a common rescaling and quantitative error remain valid.
The constraint $|\eps|=2\pi|k|/X<1$ is necessary for the present parameter
disk, but it is not by itself a uniform root-separation or convergence
estimate. A useful target is a theorem uniform when $|k|=o(X)$ with
explicit dependence on $K\lambda_k$ and the nearest singularities in the
order variable.

\subsection{Higher criticality by controlled Taylor subtraction}
The natural family yields a nondegenerate fold at its tuned critical
point, since \eqref{eq:corederivatives} never vanishes near zero.
Add finitely many adjustable analytic counterterms to force the first
several $U$ derivatives to vanish. Can one construct an equally explicit
Gamma--zeta parametrization of a cubic or higher critical locus, together
with a convergent multi-sheet chart and all-action leading coefficients?
The nondegeneracy and parameter-rank conditions must be proved before
assigning a formal multiplicity.

\subsection{Lerch and arithmetic twists}
Replace the Riemann polylogarithm by a Lerch-type series with a rational
shift or periodic weights. The local singular amplitude and regular
Laurent coefficients then involve Hurwitz zeta and generalized Stieltjes
constants. Determine precisely which periodic components preserve the
resonance lattice and which cancel its leading residue. An attractive
goal is a finite Fourier criterion for the appearance, disappearance,
or splitting of the secondary fold. No blanket extension is asserted
without checking the new local singular block.

\subsection{Vanishing analytic tails and delayed actions}
Corollary~\ref{cor:tailuniversality} shows why the leading law changes
when $c=0$. If the first nonzero tail coefficient has higher order,
classify the correct replacement for $Xe^{-X}$, locate the first new
action, and compute the sector coefficients. The center term and the
square-root discriminant may then compete on different grids; a complete
classification should retain both, rather than only the square root.

\subsection{Global zeros and inverse continuation}
Equation~\eqref{eq:physicalfold} constructs local zeros of a polylogarithm
whose order tends to one, together with a prescribed adjacent-order
value. Investigate whether these local critical points lie on distinct
global critical curves, how those curves meet the usual argument cuts,
and whether their inverse monodromies generate additional relations.
The local two-root count in this paper says nothing about all other roots
of the equation.

\subsection{Certified radii and formalization}
Produce interval-certified contour bounds for $\HH$ and $\E_n$ on
explicit disks, yielding numerical convergence radii and root enclosures
rather than floating residuals. A sensible formalization order is: the
finite coefficient algebra; the cancellation identity from a formalized
continuation formula; normal convergence and order differentiation; the
analytic root-count lemma; and the final substitutions. The supplied
manifest separates these dependencies. Existing transseries Lean modules
provide related support and inversion foundations, but this package
contains no Lean declarations and claims no new formal coverage.

\section{Conclusion}
Integer-order pole cancellation is a natural interface between the exact
polylogarithm calculus and convergent transseries inversion. Keeping that
cancellation intact reveals a second, complex resonance lattice invisible
to a fixed-order expansion. The singular core has a closed critical locus
in Gamma and zeta functions at every order. The first full-tail correction
is explicit, and the dilogarithm is exactly the order at which the
undetuned secondary scaling already has a double inverse root. Its
convergent square-root chart produces a complete leading law for every
exponential sector. These results are local and model-specific, but their
proofs, coefficient formulas, branch conventions, and replay artifacts
make them suitable for statement-level review and integration.

\begin{thebibliography}{99}
\bibitem{dlmf-poly}
NIST Digital Library of Mathematical Functions,
\S25.12, \emph{Polylogarithms}, especially Eq.~25.12.12.
\url{https://dlmf.nist.gov/25.12}. Accessed 11 October 2026.
\bibitem{dlmf-gamma}
NIST Digital Library of Mathematical Functions,
\S5.7, \emph{Series Expansions of the Gamma Function}, especially
Eq.~5.7.3. \url{https://dlmf.nist.gov/5.7}.
Accessed 11 October 2026.
\bibitem{dlmf-zeta}
NIST Digital Library of Mathematical Functions,
\S25.2, \emph{Definition and Expansions of the Riemann Zeta Function},
especially Eq.~25.2.4. \url{https://dlmf.nist.gov/25.2}.
Accessed 11 October 2026.
\bibitem{dlmf-functional}
NIST Digital Library of Mathematical Functions,
\S25.4, \emph{Reflection Formulas}, especially Eq.~25.4.2.
\url{https://dlmf.nist.gov/25.4}. Accessed 11 October 2026.
\bibitem{dlmf-w}
NIST Digital Library of Mathematical Functions,
\S4.13, \emph{Lambert $W$-Function}.
\url{https://dlmf.nist.gov/4.13}. Accessed 11 October 2026.
\bibitem{repo-poly}
ProveIt Contributors, \emph{Polylogarithms}, project index and manuscript
status. \url{https://github.com/VladimirReshetnikov/ProveIt/tree/main/Analysis/Polylogarithms}.
Inspected 11 October 2026; read scope and returned blob identifier in the
accompanying source audit.
\bibitem{repo-trans}
ProveIt Contributors, \emph{Series and Transseries}, canonical-volume,
companion, and incoming-research registry.
\url{https://github.com/VladimirReshetnikov/ProveIt/tree/main/Analysis/Transseries/docs/series-and-transseries}.
Inspected 11 October 2026, selected portions.
\bibitem{repo-hahn}
\emph{Beyond Finite-Action Folds: Critical Hahn Transseries, Stable Sector
Asymptotics, and Sharp Action Budgets}, research contribution for ProveIt,
29 September 2026. In
\path{Analysis/Transseries/docs/series-and-transseries/Critical_Hahn_Transseries_Beyond_Finite_Action_Folds/article.tex}.
The coefficient and convergence passage was inspected at commit
\texttt{9ec0f0406afe7a5475b03ff65ed91f7e283979b8}, lines 300--425.
\bibitem{repo-ledger}
ProveIt Contributors, \emph{Source Reconciliation and Mathematical
Corrections}, canonical editorial ledger, in
\path{Analysis/Polylogarithms/docs/manuscript/EDITORIAL-LEDGER.md}.
Initial reconciliation section inspected at the same pinned commit.
\bibitem{repo-bridge}
Vladimir Reshetnikov, \emph{The Polylogarithm--Polygamma Bridge: An Exact
Finite Fourier Transform, Clausen-Type Laws at Rational Arguments, and
the Layers Below the Digamma}, historical research article, 11 June 2026.
In \path{Analysis/Polylogarithms/docs/articles/polylog-polygamma-bridge.tex}.
Lines 1--160 inspected at the same pinned commit; unsupported historical
claims are superseded by the canonical ledger, not adopted here.
\end{thebibliography}
\end{document}
